\subsection{China Household Finance Survey (CHFS)}% duplicate label removed (stale file)

The China Household Finance Survey (CHFS) provides a third data source
that addresses two limitations of the PBoC and CFPS designs. The CHFS
elicits inflation expectations on a five-point scale rather than the
three-point scale used by the PBoC and CFPS, and it records household
information channels, enabling direct tests of whether the source of
price signals shapes the degree of overreaction.

The CHFS is a biennial nationally representative household survey
administered by the Survey and Research Center for China Household
Finance at Southwestern University of Finance and Economics (SWUFE)
since 2011 \citep{Gan2014}. The 2011 baseline wave sampled
approximately 8,400 households across 25 provinces and 80 counties. By
2021 the sample had expanded to over 40,000 households. The survey
collects detailed data on household assets, liabilities, income,
consumption, and financial attitudes, making it the most granular
source of household balance-sheet information available for China.

\paragraph{Inflation expectations (A4008).}
The 2011 wave asks respondents: ``Over the next year, how do you expect
commodity prices to change?'' with five response categories: rise a
lot, rise a little, no change, reduce a little, reduce a lot. This
five-point scale is richer than the three-point PBoC and CFPS measures.
The additional gradations in the tails allow me to distinguish
households who expect moderate price increases from those who expect
large ones, a distinction that maps directly onto the intensity of
diagnostic overreaction. The price-expectation question appears only in
the 2011 wave; subsequent waves dropped it in favor of questions on
housing-price expectations.

\paragraph{Information channels (A4001).}
The same wave asks households to identify their primary sources of
economic and financial information from six categories: newspapers, TV,
radio, internet, SMS alerts, and family or friends. This variable is
central to the mechanism test in Section~\ref{sec:chfs_info_channel}.
Media-intensive channels (TV, internet, newspapers) expose households
to salient price narratives that the diagnostic model predicts will
amplify overreaction. Interpersonal channels (family, friends) transmit
more diffuse and less narrative-driven signals.

\paragraph{Risk attitude (A4003).}
A self-reported risk-attitude measure on a five-point scale from very
risk-averse to very risk-loving is a placebo in the mechanism
test. If overreaction operates through salience rather than risk
preferences, the information-channel interaction should predict
forecast errors while risk attitude should not.

\paragraph{Panel structure.}
Although the expectation question is available only in 2011, the CHFS
master files track the same households through 2013, 2015, 2017, 2019,
and 2021. I link the 2011 belief measure forward to subsequent waves
using the household identifier. This panel linkage enables tests of
whether households classified as overreactors in 2011 display
systematically different income trajectories in later years.
Consumption data are not consistently available across waves, so the
forward-tracking analysis uses household income growth as the outcome.

Table~\ref{tab:chfs_2011_summary} reports summary statistics for the
2011 wave variables used in this paper.

\begin{table}[htbp]
\centering
\begin{threeparttable}
\caption{CHFS 2011 Cross-Section: Descriptive Statistics}
\label{tab:chfs_2011_summary}
{\small
\begin{tabular}{lcccccc}
\toprule
 & (1) & (2) & (3) & (4) & (5) & (6) \\
 & $N$ & Mean & Std.\ Dev. & Min & Median & Max \\
\midrule
\multicolumn{7}{l}{\textit{Panel A: Expectations and Forecast Errors}} \\
\addlinespace
Inflation expectation (1--5) & 8,333 & 4.104 & 0.837 & 1 & 4.00 & 5.000 \\
High inflation exp.\ (binary) & 8,333 & 0.841 & 0.365 & 0 & 1.00 & 1.000 \\
Forecast error proxy & 8,333 & 0.104 & 0.837 & $-3$ & 0.00 & 1.000 \\
\addlinespace
\multicolumn{7}{l}{\textit{Panel B: Information Channels}} \\
\addlinespace
High-salience channels (0--3) & 8,438 & 1.437 & 0.785 & 0 & 1.00 & 3.000 \\
Low-salience channels (0--2) & 8,438 & 0.410 & 0.561 & 0 & 0.00 & 2.000 \\
Above-median media (binary) & 8,438 & 0.383 & 0.486 & 0 & 0.00 & 1.000 \\
\addlinespace
\multicolumn{7}{l}{\textit{Panel C: Controls}} \\
\addlinespace
Risk attitude (1--5) & 8,295 & 3.847 & 1.234 & 1 & 4.00 & 5.000 \\
Log(income $+ 1$) & 8,381 & 9.768 & 2.268 & 0 & 10.24 & 14.914 \\
\bottomrule
\end{tabular}
}
\begin{tablenotes}[flushleft]
\footnotesize
\item \textit{Notes:} Data are from the 2011 China Household Finance Survey. Inflation expectation is reverse-coded so higher values indicate more inflationary expectations. Forecast error proxy equals expected (reversed ordinal) minus realized ordinal (2012 CPI $\approx 2.6\%$).
\end{tablenotes}
\end{threeparttable}
\end{table}


\paragraph{Limitations.}
Two constraints shape how the CHFS enters the analysis. Because
the expectation question appears only in 2011, I cannot construct
within-household expectation revisions across waves. All
cross-sectional tests in Section~\ref{sec:chfs_info_channel} exploit
between-household variation in a single wave rather than the
within-household panel variation available in the CFPS. The
2011 wave also predates many of the high-salience episodes that drive
the aggregate results (trade war, pork-price spike, pandemic). The CHFS
evidence therefore complements rather than replicates the time-series
and CFPS tests. Its value lies in the information-channel variation and
the richer expectation scale, not in temporal coverage.


